% ===================================================================
%  TABEL Nr. 16 — Suur kauplus. --- Basaar. Mänguasjad.
%  Traduction 2026 d'apres texto/io, controlee sur texto/fr et
%  texto/fi. Consigne : voir texto/et/10-tabel-01.tex.
%
%  DERNIER TABLEAU DE LA DERNIERE DES DIX-SEPT COLONNES, ET IL DONNE
%  DEUX FOIS LA MEME REPONSE : UNE SUR LA MENAGERIE, L'AUTRE SUR LA
%  ROSE DES VENTS.
%
%  1. LA MENAGERIE DE CARTON DE L'ALINEA 3. Le tableau 16 lituanien
%     avait vu que le chameau y est un compose, « kupranugaris », et
%     que c'est l'Evangile qui l'avait nomme. Le tableau 16 romanche
%     avait vu qu'au contraire tout y est emprunte, et que c'est le
%     JOUET qui avait nomme, en arrivant de Nuremberg avec son
%     etiquette imprimee.
%     L'ESTONIEN COMPOSE TOUT :
%       « ninasarvik »  le rhinoceros, le corne-au-nez ;
%       « jõehobu »     l'hippopotame, le cheval-de-riviere ;
%       « kilpkonn »    la tortue, la grenouille-a-bouclier ;
%       « nahkhiir »    la chauve-souris, la souris-de-cuir ;
%       « jaanalind »   l'autruche.
%     Il a repeint l'etiquette.
%
%  2. LA ROSE DES VENTS DE L'ALINEA 15. Toute l'Europe a pris les
%     quatre points cardinaux aux marins germaniques — nord, sud,
%     est, ouest, et jusqu'au russe qui garde « норд » a cote de
%     « север ». L'ESTONIEN NON :
%       « põhi »   le nord, qui veut aussi dire LE FOND ;
%       « lõuna »  le sud, qui veut aussi dire MIDI, et le repas de
%                  midi ;
%       « ida »    l'est, du lever ;
%       « lääs »   l'ouest, du coucher.
%     On s'y oriente par le soleil et par l'heure du repas, non par
%     une aiguille.
%
%  ET C'EST LA CONCLUSION DES DIX-SEPT COLONNES. Elles ont trouve,
%  l'une apres l'autre, que la forme d'un mot se lit dans le monde :
%  qui tient l'objet, qui l'a fonde, qui le distribue, par ou il est
%  venu, ou on l'apprend. Le tableau 16 lituanien disait le livre ; le
%  romanche disait le jouet.
%
%  MAIS TOUTES CES REGLES DECRIVENT CE QUI ARRIVE QUAND PERSONNE NE
%  DECIDE. Les deux dernieres colonnes de la serie sont les deux
%  seules langues qu'on ait faites — le rumantsch grischun de 1982,
%  l'estonien d'Aavik —, et ce sont les deux seules ou l'on voit ce
%  qui se passe quand quelqu'un decide : le romanche a garde ce que
%  l'usage lui donnait et n'a unifie que l'orthographe, et il suit
%  donc les regles ; l'estonien a refait les mots, et il les enfreint
%  a volonte. Il a refuse « velo » (t2), « pompier » (t11), le mot du
%  jouet (t16), et il a accepte « jootraha » (t13) et « kartul »
%  (t15) — chaque fois selon un choix, non selon une pente.
%
%  LE LIEU D'APPRENTISSAGE DECIDE DE LA FORME D'UN MOT TANT QU'IL N'Y
%  A PERSONNE POUR EN DECIDER AUTREMENT. C'est la seule chose que les
%  seize colonnes precedentes ne pouvaient pas dire, faute d'un
%  temoin ; la dix-septieme le dit.
%
%  ECART DECLARE DU BLOC t16-15-1, LE TROISIEME ET DERNIER DES TROIS.
%  Le fac-simile ido y ecrit « mi-sferi d) » sans la parenthese
%  ouvrante ; elle est retablie. UNE TRADUCTION N'EST PAS UNE
%  TRANSCRIPTION.
% ===================================================================

%%K t16-apar-1 apar
\VUsaut{4.0mm}
\VUtitre{15.0pt}{60}{TABEL Nr. 16}
\VUsaut{2.0mm}
\VUpk{13.2pt}{40}{Suur kauplus. --- Basaar.}
\VUpk{13.2pt}{40}{Mänguasjad.}
\VUsaut{3.4mm}

%%K t16-01-1 p
1. --- Uusaasta läheneb. Suured kauplused on valmistanud ette oma \VUgras{mänguasjade}
\textsuperscript{(1)} ja uusaastakinkide väljapanekud.

%%K t16-01-2 p
Ma palusin oma vanaisal viia mind basaari, et näha seda ilusat näitust, ja ta nõustus.

%%K t16-02-1 p
2. --- Kõigepealt peatusime me ilusate relvakilpide ees, milles on \VUgras{püss},
\VUgras{mundrimüts}, \VUgras{mõõk}, \VUgras{mõõgavöö} ja paar \VUgras{epolette}, või
\VUgras{kiiver}, \VUgras{kürass} \textsuperscript{(3)} ja \VUgras{püstol}
\textsuperscript{(4)}. Kõik see huvitas mind palju rohkem kui \VUgras{võlulaternad}
\textsuperscript{(2)}.

%%K t16-tit-1 sub
\VUsaut{3.0mm}
\VUpk{10.2pt}{40}{Ilusad vaatepildid.}
\VUsaut{2.6mm}

%%K t16-03-1 p
3. --- Samas osakonnas oli \VUgras{vahisõdur} \textsuperscript{(5)} oma
\VUgras{vahiputkas}; ta näis valvavat terve pappmenaaži ees. Kui palju loomi seal oli:
\VUgras{papagoi} \textsuperscript{(6)} oma õrrel, \VUgras{ninasarvik}
\textsuperscript{(7)} sarvega ninal, \VUgras{põhjapõder} \textsuperscript{(8)},
\VUgras{jaanalind} \textsuperscript{(9)}, \VUgras{ahv} \textsuperscript{(10)} pähklit
närimas, \VUgras{rebane} \textsuperscript{(11)} kana ära viimas, \VUgras{krokodill}
\textsuperscript{(12)} tohutuid lõugu ammuli ajamas, \VUgras{kaamel}
\textsuperscript{(13)}, elegantne \VUgras{hurt} \textsuperscript{(14)}, \VUgras{panter}
\textsuperscript{(15)}, siis kaks looma, keda ma ei tundnud ja kes on, nagu öeldakse,
\VUgras{nirk} \textsuperscript{(16)} ja \VUgras{hüään} \textsuperscript{(17)}. Ka seal olid
\VUgras{leopard} \textsuperscript{(18)} ja \VUgras{hunt} \textsuperscript{(21)}; päris
lähedal olid nunnud \VUgras{rotikesed} \textsuperscript{(19)} ja pisikesed
\VUgras{hiirekesed} \textsuperscript{(20)} kummist. Lõpuks lõpetasid kogu
\VUgras{jõehobu} \textsuperscript{(22)} ja küüruga \VUgras{dromedar}
\textsuperscript{(23)}. Ma oleksin sinna jäänud veel mitmeks tunniks, kui kõrval ei oleks
olnud suurepärast \VUgras{tinasõdureid} \textsuperscript{(29)} täis laagrit, mis mu
tähelepanu köitis.

%%K t16-tit-2 sub
\VUsaut{4.0mm}
\VUpk{10.2pt}{40}{Sõdurid ja sõda.}
\VUsaut{2.8mm}

%%K t16-04-1 p
4. --- Mu vanaisa, kes on \VUgras{invaliid} \textsuperscript{(24)} ja kes pidi armeest
lahkuma \VUgras{haava} tõttu, mis tõi talle \VUgras{sõjaväemedali}, jutustab mulle väga
meelsasti \VUgras{sõjaretkedest}, milles ta osales. Ta oli õnnelik seletades mulle
väikese miniatuurarmee eri liikumisi. Rühm tõelisi sõdureid, kes basaari külastasid ---
\VUgras{pioneer} \textsuperscript{(25)}, \VUgras{alpikütt} \textsuperscript{(26)} peas
\VUgras{barett}, jalaväe \VUgras{veebel} \textsuperscript{(27)} ja \VUgras{suurtükiväe
seersant} \textsuperscript{(28)} kuldse \VUgras{paelaga} varrukal --- peatus meie juures,
et seletusi kuulata.

%%K t16-05-1 p
5. --- Mu vanaisa, väga uhke, et tal on nii hiilgav kuulajaskond, improviseeris terve
lahinguplaani.

%%K t16-05-2 p
Kindlustatud linna, oma \VUgras{vallide} ja \VUgras{kraavidega}, piiras
\VUgras{vaenlase armee}, kes pärast pikka \VUgras{pommitamist} oli valmis
\VUgras{rünnakuks}. Kuid \VUgras{piiratud}, nälja sunnitud, olid proovinud
\VUgras{väljatungi} \VUgras{piirajate} \VUgras{laagri} vastu. Olles tõrjunud
\VUgras{rünnaku} visas \VUgras{lahingus} ümber \VUgras{telklaagri}, saatsid
\VUgras{võidukad} piirajad oma \VUgras{eskadronid} taga ajama \VUgras{võidetud}
ründajaid. Need läksid \VUgras{taganemisele}, kattes \VUgras{lahinguvälja}
\VUgras{surnute} ja \VUgras{haavatutega}. \VUgras{Sanitarid} viisid neid viimaseid
\VUgras{kiirabisse}, et \VUgras{kirurg} neid hooldaks.

%%K t16-05-3 p
Vaatamata mõne \VUgras{pataljoni} kangelaslikule vastupanule, kes olid moodustanud
\VUgras{kare}, oli \VUgras{kaotus} täielik, armee ülejäänud osa pöördus
\VUgras{põgenemisele}, jättes palju \VUgras{vange}. Vaenlane läks oma \VUgras{võitu}
kroonima, vallutades linna. Ehk on see sõjaretke lõpp ja pärast \VUgras{vaherahu} saab
alla kirjutada \VUgras{rahulepingule}.

%%K t16-tit-3 sub
\VUsaut{4.4mm}
\VUpk{10.2pt}{40}{Uusaastakingid.}
\VUsaut{2.8mm}

%%K t16-06-1 p
6. --- Noored sõdurid, kes naersid, kuulates veterani maalilist jutustust, küsisid
temalt veel mõned seletused. Mina käisin ümber müügiletti, et vaadata ilusat \VUgras{ammu}
\textsuperscript{(32)} ja \VUgras{vibu} \textsuperscript{(33)} oma \VUgras{noolega}
\textsuperscript{(31)}. Seal leidsin ma teise loomade kogu: kaks \VUgras{laululindu}
\textsuperscript{(34)} --- automaatne \VUgras{ööbik} ja \VUgras{põõsalind}, kes laulsid
täiest kõrist --- ja rea kummalisi loomi: \VUgras{nahkhiir} \textsuperscript{(35)},
\VUgras{boa} \textsuperscript{(36)}, \VUgras{kilpkonn} \textsuperscript{(37)},
\VUgras{hüljes} \textsuperscript{(38)}, \VUgras{vaal} \textsuperscript{(39)} ja
täispuhutav \VUgras{haikala} \textsuperscript{(40)}.

%%K t16-07-1 p
7. --- Ma ei peatunud kaua neid vaatlemas, sest ma nägin seda, millest ma unistasin:
suurepärast \VUgras{marionettteatrit} \textsuperscript{(41)} imeliste
\VUgras{dekoratsioonide} ja võluva punase \VUgras{eesriidega}. \VUgras{Stseenil} näisid
kaks näitlejat mängivat \VUgras{rolli} väga huvitavas \VUgras{näidendis}, sest väikesed
papist \VUgras{pealtvaatajad}, kes olid seatud loožidesse või istusid
\VUgras{tugitoolidel} ja \VUgras{parteris} \VUgras{orkestrijuhi} taga, aplodeerisid
elavalt.

%%K t16-07-2 p
Kahjuks maksis see teater väga kallilt.

%%K t16-08-1 p
8. --- Pealegi oli mänguasju raske valida, sest vitriinidesse oli kuhjatud igasuguseid
mänguasju: \VUgras{Arlekiin} \textsuperscript{(42)}, \VUgras{Pulcinella}
\textsuperscript{(43)}, \VUgras{Pierrot} \textsuperscript{(44)}, \VUgras{kõrvukrätsid}
\textsuperscript{(45)} tiibu lehvitamas, vilistavad \VUgras{mustad rästad}
\textsuperscript{(46)}, haukuvad \VUgras{puudlid}, \VUgras{fonograaf}
\textsuperscript{(47)} ja \VUgras{kannatusemängud}. Üks neist mänguasjadest lõbustas mind
väga: see kujutas \VUgras{merelahingut} \textsuperscript{(48)}, milles \VUgras{laeva}
ründasid \VUgras{piraadid} maa lähedal, mis oli istutatud täis \VUgras{kookospalme}.

%%K t16-noto-1 noto
\VUnotes{143.2mm}{%
\textit{(*) Vaata tabelit nr.~6.}%
}

%%K t16-09-1 p
9. --- Ma võisin väga kergesti vaadelda kõiki neid imesid, sest kaupluses oli vähe
inimesi, vaevalt mõned uitajad, \VUgras{draguun} \textsuperscript{(49)} ja
\VUgras{inkassaator} \textsuperscript{(50)}, kes esitas \VUgras{vekslit} peamises
\VUgras{kassas} \textsuperscript{(51)}.

%%K t16-10-1 p
10. --- Ma küsisin oma vanaisalt, milleks on raamatud, mis on pandud riiulitele
\VUgras{kassapidaja} taha; ta vastas mulle, et need on
\VUgras{raamatupidamisraamatud}.

%%K t16-tit-4 sub
\VUsaut{3.6mm}
\VUpk{10.2pt}{40}{Poe ümber vahtimine.}
\VUsaut{2.8mm}

%%K t16-11-1 p
11. --- Mänguasjaosakonnast lahkudes läksime me sadulsepatöökotta, et heita pilk
\VUgras{sadulatele} \textsuperscript{(53)}, \VUgras{jalustele} \textsuperscript{(54)},
\VUgras{valjastele} \textsuperscript{(55)}, \VUgras{suulistele} \textsuperscript{(56)} ja
\VUgras{piitsadele}. Sellal kui \VUgras{osakonnajuhataja} \textsuperscript{(57)} andis mu
vanaisale mõned teated, läksin ma vaatama \VUgras{portselani} ja \VUgras{kristalli}
\textsuperscript{(58)} väljapanekut, kus on ilusad \VUgras{salatikausid} ja peened
\VUgras{teekannud} \textsuperscript{(59)}.

%%K t16-12-1 p
12. --- Et minna teisele korrusele, läksime me \VUgras{korsettide}
\textsuperscript{(60)} vitriini tagant, sukakaupluse kõrvalt, kus olid välja pandud väga
ilusad \VUgras{aluspüksid} \textsuperscript{(63)}. Lähedal laskis \VUgras{müüjanna}
\textsuperscript{(61)} \VUgras{ostjal} \textsuperscript{(62)} valida ilusaid siidiseid
\VUgras{kangaid} \textsuperscript{(64)}.

%%K t16-12-2 p
Trepist üles minnes märkasin ma, et kauplus on kaunistatud jaapani \VUgras{laternate}
\textsuperscript{(65)}, \VUgras{päikesevarjude} \textsuperscript{(66)} ja hiiglaslike
\VUgras{palmidega} \textsuperscript{(70)}.

%%K t16-13-1 p
13. --- Teisel korrusel ei olnud palju näha; ainult mööbel (*), \VUgras{seifid}
\textsuperscript{(67)}, \VUgras{kummutid} \textsuperscript{(68)}, \VUgras{lastevankrid}
\textsuperscript{(69)}. Kuid kaupluse üldvaade oli väga \VUgras{lõbus}.

%%K t16-14-1 p
14. --- Meie jalge ees nägin ma muusikariistu: \VUgras{harfid} \textsuperscript{(71)},
\VUgras{kitarrid} \textsuperscript{(72)}, \VUgras{mandoliinid} \textsuperscript{(73)},
\VUgras{ventiilkornetid} \textsuperscript{(74)}, \VUgras{klarnetid}
\textsuperscript{(75)}, viiulid oma \VUgras{poognaga} \textsuperscript{(76)} ja isegi
\VUgras{suur trumm} \textsuperscript{(77)}. Veidi kaugemal, paberiosakonna juures, tingis
\VUgras{üliõpilane} \textsuperscript{(78)} \VUgras{müüjaga} \textsuperscript{(79)}
vihikumapi pärast. Nende kõrval eristasin ma ilusat \VUgras{aabitsat}
\textsuperscript{(80)}.

%%K t16-14-2 p
Kaupluse keskel oli püstitatud kõrge \VUgras{jõulupuu} \textsuperscript{(81)}, mis oli
üleni varustatud küünalde, nipsasjade ja igasuguste mänguasjadega: \VUgras{puuhobused}
\textsuperscript{(82)}, \VUgras{nukud} \textsuperscript{(83)} jne.

%%K t16-14-3 p
\VUgras{Parfümeeria} \textsuperscript{(84)} oma \VUgras{hambavete} ja mitmesuguste
\VUgras{parfüümidega} levitas väga head lõhna, mis meieni tõusis.

%%K t16-tit-5 sub
\VUsaut{3.4mm}
\VUpk{10.2pt}{40}{Me ostame midagi.}
\VUsaut{2.8mm}

%%K t16-15-1 p
15. --- Kuna mu vanaisa hakkas aega liiga pikaks pidama, läksime me suurest trepist
alla. Ta peatus veidi \VUgras{astronoomiatabeli} \textsuperscript{(85)} ees, mis kujutas,
nagu ta mulle ütles, \VUgras{komeete} \textsuperscript{(a)}, \VUgras{kuu- ja
päikesevarjutusi} \textsuperscript{(b)}, tuuleroosi \VUgras{nelja ilmakaarega}
\textsuperscript{(c)} \VUgras{(põhi, lõuna, ida ja lääs)}, kahe \VUgras{poolkera}
\textsuperscript{(d)} kaarti, \VUgras{planeete} \textsuperscript{(e)} ja
\VUgras{päikesesüsteemi} \textsuperscript{(f)}. Ma ei saanud kõigest sellest midagi aru,
ja mind huvitas palju rohkem paelutatud livree \VUgras{väljakandjal}
\textsuperscript{(86)}, kes möödus, ja ilusad \VUgras{lehvikud} \textsuperscript{(87)},
mis olid mu kõrval, \VUgras{rahakottide} \textsuperscript{(88)} ja \VUgras{portfellide}
\textsuperscript{(89)} juures.

%%K t16-16-1 p
16. --- Ma imetlesin ka väljapanekul \VUgras{sulgi} \textsuperscript{(90)} ja
\VUgras{vööpandlaid} \textsuperscript{(91)}, ilusaid \VUgras{kunstlilli}
\textsuperscript{(94)}, mis olid graatsiliselt korvidesse seatud. \VUgras{Kannikesed},
\VUgras{levkoid}, \VUgras{hüatsindid} \textsuperscript{(95)}, \VUgras{pelargoonid}
\textsuperscript{(96)}, \VUgras{liiliad} \textsuperscript{(97)} ja
\VUgras{kapuutsililled} \textsuperscript{(98)} olid väga hästi järele tehtud.

%%K t16-tit-6 sub
\VUsaut{4.4mm}
\VUpk{10.2pt}{40}{Vanaisa kingitus.}
\VUsaut{2.8mm}

%%K t16-17-1 p
17. --- Ma palusin oma vanaisal osta mulle üks \VUgras{hambaharjadest}
\textsuperscript{(92)}, mis olid karbis \VUgras{juuksenõelte} \textsuperscript{(93)}
kõrval. Ta nõustus ja ma läksin ise oma ostu kassas maksma, mille juures valmistati
ette tähtsat \VUgras{saadetist} \textsuperscript{(100)} ühele \VUgras{kliendile}
\textsuperscript{(99)}.

%%K t16-18-1 p
18. --- Järsku tekkis kerge segadus inimeste seas, kes olid peatunud kunstiliste
\VUgras{pronksesemete} \textsuperscript{(101)} juures. \VUgras{Varas}
\textsuperscript{(102)} oli äsja teolt tabatud \VUgras{inspektori} \textsuperscript{(103)}
poolt, kui ta püüdis kedagi paljaks röövida. Hoolitseti selle eest, et politseinik
tuuakse tema \VUgras{vahistamiseks}, ja just kui me välja läksime, viidi kurjategija
vanglasse.

\VUsaut{6.0mm}
\VUfilet{22mm}
